\sporkSubsection{Kernel Structure}{sec:KernelStructure}

We open a CUDA scope using a \texttt{with CudaDeviceFunction} statement.
The statement takes a \texttt{clusterDim} parameter,
specifying the number of CTAs (thread blocks) per cluster (must be 1 for Ampere and older architectures),
as well as a list of warp configurations for warp specialization (Section~\ref{sec:WarpSpecialization}).
Figure~\ref{fig:KernelStructure} shows the overall structure of a kernel:
we require the body of the kernel to consist of a single nest of \texttt{cuda\_tasks} loops;
the inner-most \texttt{cuda\_tasks} loop body comprises a \emph{device task},
which can contain \texttt{seq} and \texttt{cuda\_threads} loops and warp specialization statements.
A single execution of a device task is assigned to one CUDA cluster.
The compiler implements a \emph{persistent kernel} strategy, where multiple device tasks may be executed by a single cluster.

\begin{figure}[h]
\sporkCodeRule
\centering
\begin{varwidth}{\linewidth}
\sporkListing{kernel_structure}
\end{varwidth}
\caption{Kernel structure.}
\label{fig:KernelStructure}
\sporkCodeRule
\end{figure}


Each thread is identified with a $\sporktid$ (thread ID), consisting of a \emph{task ID} component and \emph{local thread index} component.

% I'm just info-dumping here but I'm catastrophically out of time and it breaks my heart, it really does.

The task ID is simply the tuple of current \texttt{cuda\_tasks}-loop iterator values.
This abstracts over the specific mapping between device tasks and hardware clusters.
We currently only implement a static schedule where device tasks are assigned round-robin to CUDA clusters,
but future work could swap this with a selection of strategies (e.g. Morton swizzling)
without affecting downstream analysis.

The local thread index identifies the specific thread within the (abstracted) cluster.
Threads are numbered lexicographically by CTA rank (within the cluster), then by thread index in the CTA.
Hence, the local thread indices range from 0 to $\texttt{clusterDim}\times\texttt{blockDim}-1$;
\texttt{blockDim} being the number of threads per CTA.
The \texttt{cuda\_threads} loops and warp specialization statements decide the local thread indices given by a scope's $\sporktidMap$;
a \texttt{cuda\_threads(0, $\sporkz$)} loop effectively takes its executing thread collective and splits it into $\sporkz$-many disjoint subsets
(possibly with threads left-over),
and assigns 1 each to each loop iteration without any implied fork-join semantics.
